% Chapter 3: Core Pipeline Architecture & High-Throughput Serving

\begin{frame}[plain]
    \vfill
    \centering
    {\footnotesize \textbf{\textcolor{deepblue}{\textsc{TaxonVision-MLOps Architecture}}}} \\[0.35cm]
    {\large \textbf{\textcolor{signalorange}{CHAPTER 03}}} \\[0.25cm]
    {\huge \textbf{Pipeline Architecture \& Fast-Path Serving}} \\[0.35cm]
    {\normalsize \textcolor{black!70}{Decoupled Representation, High-Throughput ONNX \& Latency Budgets}} \\[0.5cm]
    \textcolor{deepblue}{\rule{0.55\linewidth}{1.2pt}} \\[0.5cm]
    \statusdone
    \vfill
    \script{In Chapter 3, we present our decoupled pipeline architecture: separating heavy self-supervised feature extraction from lightweight linear heads and sub-15ms production serving.}
    \qa{What is the core breakthrough of this chapter?}{Decoupling static representation from downstream classification, achieving 10-second training cycles and zero-allocation high-throughput inference serving.}
\end{frame}

\begin{frame}{Core Architecture}{Decoupling Representation from Reasoning}
    \begin{block}{The Decoupling Principle}
        \begin{itemize}\small
            \item \textbf{The Monolithic Trap:} Junior practitioners often attempt to train an end-to-end convolutional or vision transformer network that directly outputs species logits.
            \item \textbf{The Retraining Bottleneck:} When new species are discovered or taxonomic revisions occur, an entire multi-gigabyte backbone must be painfully retrained.
            \item \textbf{The Decoupled Solution:} TaxonVision cleanly decouples static visual representation (offline feature extraction) from dynamic head training and high-throughput inference serving.
            \item \textbf{Agile Outcomes:} This separation enables agile, localized updates, rapid hyperparameter sweeps, and sub-15ms classification without full retraining loops.
        \end{itemize}
    \end{block}

    \vspace{0.15cm}
    \statusdone

    \script{This slide represents the core architectural breakthrough of our system: decoupling representation from classification. In naïve computer vision, you take an image, pass it through an end-to-end network, and output 1,000 class probabilities. But in biodiversity, labels and taxonomy evolve constantly. By decoupling representation from downstream heads, our heavy foundation model stays entirely frozen, while lightweight linear probes train in seconds and deploy instantly.}
    \qa{What is a foundation model backbone?}{A vision transformer (BioCLIP, DINOv2) pre-trained on vast imagery using self-supervised learning, yielding rich morphological representations without task-specific retraining.}
    \qa{How fast does head training become with frozen embeddings?}{Because feature embeddings are pre-extracted and cached in Parquet/NumPy arrays, training a linear probe takes only 10 to 30 seconds rather than hours of multi-GPU compute.}
\end{frame}

\begin{frame}{Pipeline Architecture}{Offline Representation \& Linear Probe Training}
    \begin{block}{Offline Batch Representation Engine}
        \begin{itemize}\small
            \item \textbf{S3 Batch Streaming:} High-resolution raw images stored in remote S3 are processed in zero-copy batch streams.
            \item \textbf{Frozen Backbone Inference:} Images pass through frozen backbones (DINOv2, BioCLIP) to generate dense feature embeddings (768 or 1024 dimensions).
            \item \textbf{Serialization:} These embeddings are serialized into content-addressable storage (Parquet/DVC).
            \item \textbf{Linear Probing:} Downstream, lightweight linear classifiers are trained using Class-Balanced Loss and Bayesian hyperparameter sweeps (Optuna).
            \item \textbf{ONNX Export:} The final heads are compiled into ONNX-quantized models ready for high-throughput serving.
        \end{itemize}
    \end{block}

    \vspace{0.15cm}
    \statusdone

    \script{This is our offline training engine. Instead of feeding raw images directly into training loops every epoch, we pass them through our frozen backbone once and store the resulting numerical embeddings on disk. Once those embeddings are saved, training a new classification head takes literally ten seconds. We can run 50 hyperparameter optimization trials with Optuna in five minutes on a single GPU—something that would take days with full end-to-end fine-tuning.}
    \qa{Does training a linear head hurt accuracy?}{Self-supervised backbones produce features where linear probes achieve within 1--2\% of full fine-tuning, while saving orders of magnitude in compute and preventing catastrophic forgetting.}
    \qa{What is Class-Balanced Loss?}{A loss weighting technique correcting for long-tail biological distributions, ensuring rare species are not ignored by the model.}
\end{frame}

\begin{frame}{Inference Serving}{High-Throughput Serving \& Zero-Allocation Fast-Path}
    \begin{block}{Production Serving Invariants}
        \begin{itemize}\small
            \item \textbf{Zero-Allocation Invariant:} ONNX Runtime inference hot paths pre-allocate execution buffers, completely eliminating dynamic memory allocations and Python object boxing.
            \item \textbf{Strict Latency Budget:} Sub-15ms p95 latency target on bare-metal hardware, maintaining high throughput under concurrent HTTP workloads.
            \item \textbf{INT8 Dynamic Quantization:} Quantizes dense linear probe weights to 8-bit integers, reducing cache footprint by $4\times$ while retaining $\ge 99\%$ accuracy.
            \item \textbf{Conformal Safety Interlock:} Prediction logits are evaluated simultaneously against pre-computed conformal non-conformity thresholds before response delivery.
        \end{itemize}
    \end{block}

    \vspace{0.15cm}
    \statusdone

    \script{In production APIs, speed and determinism are everything. We implement a strict zero-allocation fast-path serving boundary: incoming images are processed through pre-allocated tensor buffers in ONNX Runtime. The quantized linear probe executes in under 15 milliseconds, immediately paired with conformal prediction intervals to guarantee statistical coverage.}
    \qa{What is ONNX Runtime and why do we use it instead of native PyTorch?}{ONNX Runtime optimizes the computation graph, fuses operations, and performs strict memory reuse. It eliminates Python interpreter overhead, resulting in 2x to 4x lower latency and rigorously enforces zero memory allocations on the hot inference path.}
    \qa{What is the purpose of INT8 quantization here?}{It reduces memory bandwidth pressure and allows classification heads to fit entirely within L2/L3 CPU caches or fast GPU registers, eliminating cache thrashing.}
\end{frame}

\begin{frame}{Architectural Topology}{Offline Training to Online Serving Boundary}
    \begin{center}
    \vspace{-0.2cm}
    \providecolor{TaxonPrimary}{RGB}{34, 139, 34} % plantgreen
    \providecolor{TaxonSecondary}{RGB}{220, 120, 0} % signalorange
    \providecolor{TaxonDark}{RGB}{0, 51, 102} % deepblue
    \providecolor{TaxonAccent}{RGB}{180, 30, 30} % herbired
    \resizebox{0.9\textwidth}{!}{
    \begin{tikzpicture}[
        node distance=1.5cm and 2.0cm,
        box/.style={rectangle, draw=TaxonDark, fill=TaxonDark!10, thick, rounded corners, minimum width=2.5cm, minimum height=1cm, align=center, font=\sffamily\small},
        fastapi/.style={rectangle, draw=TaxonPrimary, fill=TaxonPrimary!15, thick, rounded corners, minimum width=4cm, minimum height=1.2cm, align=center, font=\sffamily\bfseries},
        arrow/.style={->, >=stealth, thick, draw=TaxonDark!70},
        dashedarrow/.style={->, >=stealth, thick, dashed, draw=TaxonSecondary}
    ]

    % Row 1: Offline Training Pipeline
    \node[box] (s3) {S3 Images\\(Zero-Copy)};
    \node[box, right=of s3] (vit) {Frozen ViT\\(DINOv2 / BioCLIP)};
    \node[box, right=of vit] (embed) {Embeddings\\(Parquet / DVC)};
    \node[box, right=of embed] (linear) {Linear Head\\(Optuna / CB-Loss)};
    \node[box, right=of linear] (onnx) {ONNX Export\\(INT8 Quantized)};

    % Row 2: Online Serving Path
    \node[box, below=1.8cm of s3] (ingress) {Client Image\\Upload};
    \node[box, right=of ingress] (prep) {Zero-Copy\\Preprocess};
    \node[box, right=of prep] (session) {ONNX Runtime\\Inference Engine};
    \node[box, right=of session] (conformal) {Conformal\\Safety Gate};
    \node[fastapi, right=of conformal] (api) {FastAPI Response\\p95 $< 15$ ms};

    % Row 1 Arrows
    \draw[arrow] (s3) -- (vit);
    \draw[arrow] (vit) -- (embed);
    \draw[arrow] (embed) -- (linear);
    \draw[arrow] (linear) -- (onnx);

    % Row 2 Arrows
    \draw[arrow] (ingress) -- (prep);
    \draw[arrow] (prep) -- (session);
    \draw[arrow] (session) -- (conformal);
    \draw[arrow] (conformal) -- (api);

    % Connection from Artifact to Serving Engine
    \draw[dashedarrow] (onnx) -- node[right, font=\sffamily\scriptsize, text=TaxonSecondary] {Artifact Release} (session);

    % Grouping boxes
    \begin{scope}[on background layer]
        \node[draw=TaxonDark!30, fill=TaxonDark!5, dashed, rounded corners, fit=(s3)(onnx), inner sep=0.3cm, label={[font=\sffamily\bfseries, text=TaxonDark]above:Offline Training Pipeline (Batch)}] {};
        \node[draw=TaxonPrimary!40, fill=TaxonPrimary!5, dashed, rounded corners, fit=(ingress)(api), inner sep=0.3cm, label={[font=\sffamily\bfseries, text=TaxonPrimary]below:Online Fast-Path Serving (Real-Time)}] {};
    \end{scope}

    \end{tikzpicture}
    }
    \end{center}
    \script{This diagram visualizes our decoupled topology. The top row represents offline training: zero-copy S3 streaming feeds frozen ViT backbones, producing Parquet embedding manifests that train linear heads and compile into ONNX artifacts. The bottom row represents the online inference serving path: requests enter FastAPI, execute in pre-allocated ONNX Runtime sessions, pass through conformal safety gates, and respond in under 15 milliseconds.}
    \qa{How does the serving engine update when a new head is trained?}{The compiled ONNX model artifact is tagged and versioned in the model registry. The FastAPI serving pod dynamically mounts the new artifact without interrupting active worker threads.}
\end{frame}

\begin{frame}{Decision Flowchart}{Fast-Path vs. Epistemic Triage Routing}
    \begin{center}
    \vspace{-0.2cm}
    \providecolor{TaxonPrimary}{RGB}{34, 139, 34}
    \providecolor{TaxonSecondary}{RGB}{220, 120, 0}
    \providecolor{TaxonDark}{RGB}{0, 51, 102}
    \resizebox{0.95\linewidth}{!}{
    \begin{tikzpicture}[
        node distance=0.8cm and 1.2cm,
        start/.style={draw=TaxonDark, fill=TaxonDark!20, ellipse, minimum width=2.2cm, minimum height=0.9cm, align=center, font=\sffamily\bfseries\scriptsize},
        process/.style={draw=TaxonDark, fill=TaxonDark!10, rectangle, rounded corners, minimum width=2.6cm, minimum height=0.9cm, align=center, font=\sffamily\scriptsize},
        decision/.style={draw=TaxonSecondary, fill=TaxonSecondary!15, diamond, aspect=1.8, minimum width=2.5cm, minimum height=1.1cm, align=center, font=\sffamily\scriptsize},
        endpoint/.style={draw=TaxonPrimary, fill=TaxonPrimary!20, rectangle, rounded corners, thick, minimum width=2.4cm, minimum height=0.9cm, align=center, font=\sffamily\bfseries\scriptsize},
        oodpoint/.style={draw=herbired, fill=herbired!15, rectangle, rounded corners, thick, minimum width=2.4cm, minimum height=0.9cm, align=center, font=\sffamily\bfseries\scriptsize},
        arrow/.style={->, >=stealth, thick, draw=TaxonDark!80}
    ]

    % Horizontal Pipeline
    \node[start] (req) {Inference\\Request};
    \node[process, right=of req] (onnx) {ONNX Runtime\\(Linear Head)};
    \node[process, right=of onnx] (conformal) {Split Conformal\\Prediction Set};
    \node[decision, right=of conformal] (size) {Set Size\\$|C(x)|$?};

    % Branch Outcomes
    \node[endpoint, above right=0.3cm and 1.2cm of size] (fast) {\textbf{Fast-Path Return}\\($|C|=1$, $<15$ ms)};
    \node[process, right=1.2cm of size] (ambig) {Ambiguity Triage Queue\\(Defer Deep Retrieval)};
    \node[endpoint, right=0.8cm of ambig] (enriched) {\textbf{Triage Output}\\(Multi-Taxa Review)};
    \node[oodpoint, below right=0.3cm and 1.2cm of size] (ood) {\textbf{Human Triage}\\($|C|=0$, Epistemic OOD)};

    \draw[arrow] (req) -- (onnx);
    \draw[arrow] (onnx) -- (conformal);
    \draw[arrow] (conformal) -- (size);

    \draw[arrow] (size) |- node[above, font=\sffamily\tiny, pos=0.7] {$|C|=1$ (Singleton)} (fast);
    \draw[arrow] (size) -- node[above, font=\sffamily\tiny] {$|C|>1$ (Ambiguous)} (ambig);
    \draw[arrow] (ambig) -- (enriched);
    \draw[arrow] (size) |- node[below, font=\sffamily\tiny, pos=0.7] {$|C|=0$ (Epistemic OOD)} (ood);

    \end{tikzpicture}
    }
    \end{center}
    \script{Here we zoom into the FastAPI serving decision logic. Every request runs through ONNX Runtime and the conformal calibration layer. If the conformal set size is exactly one, the model is mathematically confident, and we immediately return via the Fast-Path in under 15 milliseconds. If the set contains multiple candidate taxa, the observation is flagged as ambiguous and routed to our triage queue for deeper analysis. If the set is empty, it is flagged as an epistemic outlier.}
    \qa{What happens when $|C(x)| > 1$?}{Ambiguous observations route to our review queue where multi-taxa ties are flagged. Later in Chapter 10, we will see how our Multimodal Knowledge Graph and Vision-RAG architecture automates deep differential diagnosis for these exact cases.}
    \qa{What happens if the conformal set is empty ($|C(x)| = 0$)?}{An empty set implies the image is an epistemic or out-of-distribution outlier. In this case, the system safely routes the observation to human review with an unclassified referral flag.}
\end{frame}
